\documentclass[11pt,a4paper]{article}

% Pakete
\usepackage[utf8]{inputenc}
\usepackage[T1]{fontenc}
\usepackage[ngerman]{babel}
\usepackage[left=18mm,right=18mm,top=18mm,bottom=18mm]{geometry}
\usepackage{helvet}
\renewcommand{\familydefault}{\sfdefault}
\usepackage{amsmath,amssymb}
\usepackage{graphicx}
\usepackage{tikz}
\usetikzlibrary{arrows.meta,positioning}
\usepackage{tcolorbox}
\tcbuselibrary{skins,breakable}
\usepackage{enumitem}
\usepackage{tabularx}
\usepackage{colortbl}
\usepackage{siunitx}
\sisetup{locale=DE, per-mode=fraction}
\usepackage{qrcode}

% Fachfarbe Physik
\definecolor{physik}{HTML}{2c5aa0}
\definecolor{physiklight}{HTML}{e3f2fd}
\definecolor{loesung}{HTML}{c62828}

% Box-Stile
\tcbset{
    merkbox/.style={
        colback=physiklight,
        colframe=physik,
        fonttitle=\bfseries,
        title=#1,
        boxrule=1pt,
        arc=2pt,
        left=5pt,
        right=5pt,
        top=5pt,
        bottom=5pt
    },
    konstantenbox/.style={
        colback=white,
        colframe=physik,
        boxrule=1pt,
        arc=2pt,
        left=6pt,
        right=6pt,
        top=4pt,
        bottom=4pt
    }
}

% Lösung-Befehl
\newcommand{\lsg}[1]{\textcolor{loesung}{\textbf{#1}}}

% Kopfzeile
\usepackage{fancyhdr}
\pagestyle{fancy}
\fancyhf{}
\lhead{\small Physik 12 GK}
\chead{\small \textcolor{loesung}{MUSTERLÖSUNG}}
\rhead{\small Quantenphysik}
\renewcommand{\headrulewidth}{0.4pt}

% Keine Einrückung
\setlength{\parindent}{0pt}
\setlength{\parskip}{6pt}

\begin{document}

% Titel mit QR-Codes
\begin{minipage}[t]{0.62\textwidth}
\vspace{0pt}
{\Large\textbf{Photoeffekt Vertiefung + De-Broglie-Wellenlänge}}\\[0.3em]
{\normalsize Fr 16.01.2026 \hfill \textcolor{loesung}{MUSTERLÖSUNG}}
\end{minipage}
\hfill
\begin{minipage}[t]{0.35\textwidth}
\vspace{0pt}
\raggedleft
\begin{tabular}{@{}c@{\hspace{8pt}}c@{}}
\qrcode[height=1.3cm]{https://devbox42.github.io/sciencesim/physik/kl12-GK/01-photoeffekt/03-DS-de-broglie/SIM-03-de-broglie.html} &
\qrcode[height=1.3cm]{https://devbox42.github.io/sciencesim/physik/kl12-GK/01-photoeffekt/03-DS-de-broglie/LP-03-de-broglie.html} \\
{\scriptsize Simulation} & {\scriptsize Lernpfad}
\end{tabular}
\end{minipage}

\vspace{0.3em}
\hrule height 1pt
\vspace{0.5em}

% ============================================================
% KONSTANTEN (kompakt)
% ============================================================

\begin{tcolorbox}[konstantenbox]
\small
\textbf{Konstanten:}\quad
$h = \SI{6,63e-34}{J.s} = \SI{4,14e-15}{eV.s}$ \hfill
$e = \SI{1,6e-19}{C}$ \hfill
$m_e = \SI{9,1e-31}{kg}$ \hfill
$c = \SI{3e8}{m/s}$
\end{tcolorbox}

\vspace{0.3em}

% ============================================================
% BLOCK 1: PHOTOEFFEKT
% ============================================================

{\large\textbf{Block 1: Photoeffekt}}

\vspace{0.3em}

\textbf{Aufgabe 1: Photoeffekt an Cäsium} \hfill \textbf{(11 BE)}

In einem Physiklabor wird der Photoeffekt an einer Cäsium-Kathode untersucht. Die Austrittsarbeit von Cäsium beträgt $W_A = \SI{1,9}{eV}$. Das Diagramm zeigt die Messergebnisse.

\vspace{0.3em}

\begin{center}
\begin{tikzpicture}[xscale=0.9, yscale=0.9]
    % Gitter
    \draw[gray!30, very thin, xstep=1, ystep=0.5] (0,-2) grid (9,2);

    % Achsen
    \draw[->, thick] (0,-2) -- (0,2.3) node[above] {$E_{\text{kin}}$ (eV)};
    \draw[->, thick] (0,0) -- (9.5,0) node[right] {$f$ ($10^{14}$ Hz)};

    % Achsenbeschriftung x
    \foreach \x in {1,2,3,4,5,6,7,8,9} {
        \draw (\x,0.1) -- (\x,-0.1) node[below] {\small \x};
    }
    % Achsenbeschriftung y
    \foreach \y in {-2,-1,1,2} {
        \draw (0.1,\y) -- (-0.1,\y) node[left] {\small \y};
    }
    \node[below left] at (0,0) {\small 0};

    % Einstein-Gerade für Cäsium: E_kin = 0,414·f - 1,9
    % f=0: E=-1,9; f=4,59: E=0; f=9: E=1,83
    \draw[thick, physik] (0,-1.9) -- (9,1.83);

    % Schnittpunkt mit x-Achse markieren (f_G = 4,59)
    \fill[loesung] (4.59,0) circle (3pt);
    \node[loesung, above right, font=\small\bfseries] at (4.59,0.1) {$f_G \approx 4{,}6$};

    % y-Achsenabschnitt markieren
    \fill[loesung] (0,-1.9) circle (3pt);
    \node[loesung, right, font=\small\bfseries] at (0.1,-1.9) {$-W_A = -1{,}9$};
\end{tikzpicture}
\end{center}

\begin{enumerate}[label=\alph*), leftmargin=*, nosep]
\item \textbf{Ermittle} aus dem Diagramm die Grenzfrequenz $f_G$ von Cäsium. \hfill (2 BE)

\vspace{0.3em}
$f_G = $ \lsg{4,6} $\cdot 10^{14}$ Hz

\lsg{Ablesen: x-Achsenabschnitt bei $E_{\text{kin}} = 0$}

\vspace{0.5em}

\item \textbf{Berechne} die kinetische Energie eines Photoelektrons bei Bestrahlung mit UV-Licht der Frequenz $f = \SI{7e14}{Hz}$. \hfill (3 BE)

\vspace{0.3em}
$E_{\text{kin}} = h \cdot f - W_A$

\vspace{0.3em}
$E_{\text{kin}} = $ \lsg{$\SI{4,14e-15}{eV.s}$} $\cdot$ \lsg{$\SI{7e14}{Hz}$} $-$ \lsg{$\SI{1,9}{eV}$}

\vspace{0.3em}
$E_{\text{kin}} = $ \lsg{$\SI{2,9}{eV}$} $-$ \lsg{$\SI{1,9}{eV}$} $=$ \lsg{$\SI{1,0}{eV}$}

\vspace{0.5em}

\item \textbf{Erkläre}, warum der y-Achsenabschnitt der Geraden negativ ist. Gib seine physikalische Bedeutung an. \hfill (3 BE)

\vspace{0.3em}
\lsg{Der y-Achsenabschnitt entspricht $E_{\text{kin}}$ bei $f = 0$. Da bei $f = 0$ keine Photonenenergie vorhanden ist, würde das Elektron die negative Energie $-W_A$ haben (,,Energieschuld''). Physikalisch bedeutet der Betrag die Austrittsarbeit $W_A = \SI{1,9}{eV}$, die mindestens aufgebracht werden muss.}

\vspace{0.5em}

\item \textbf{Begründe}, warum eine Erhöhung der Lichtintensität bei $f = \SI{3e14}{Hz}$ keinen Photostrom erzeugt. \hfill (3 BE)

\vspace{0.3em}
\lsg{Die Frequenz $f = \SI{3e14}{Hz}$ liegt unterhalb der Grenzfrequenz $f_G = \SI{4,6e14}{Hz}$. Die Energie eines einzelnen Photons ($E = h \cdot f = \SI{1,24}{eV}$) ist kleiner als die Austrittsarbeit $W_A = \SI{1,9}{eV}$. Eine höhere Intensität bedeutet nur mehr Photonen, aber jedes einzelne hat zu wenig Energie, um ein Elektron herauszulösen.}

\end{enumerate}

% ============================================================
% KASTEN 1: MERKSATZ
% ============================================================

\begin{tcolorbox}[merkbox={Kasten 1: Merksatz Photoeffekt}]
Beim Photoeffekt gilt: Ein \lsg{Photon} überträgt seine gesamte Energie auf \lsg{ein} Elektron.

\vspace{0.3em}
Unterhalb der \lsg{Grenzfrequenz} ($f < f_G$) findet kein Photoeffekt statt -- unabhängig von der \lsg{Intensität} des Lichts.
\end{tcolorbox}

\newpage

% ============================================================
% BLOCK 2: DE-BROGLIE
% ============================================================

{\large\textbf{Block 2: De-Broglie-Wellenlänge}}

\vspace{0.3em}

% ============================================================
% KASTEN 2: DE-BROGLIE
% ============================================================

\begin{tcolorbox}[merkbox={Kasten 2: De-Broglie-Hypothese (1924)}]
\textbf{Symmetrieargument:} Wenn Licht Teilcheneigenschaften hat (Photonen), haben dann Teilchen auch Welleneigenschaften?

\vspace{0.5em}
\textbf{De Broglies Antwort:} Jedes bewegte Teilchen hat eine Wellenlänge:

\begin{center}
\fbox{\rule{0pt}{1.2em}\hspace{0.5em}$\displaystyle \lambda = \frac{h}{p} = \frac{h}{m \cdot v}$\hspace{0.5em}}
\end{center}

\vspace{0.3em}
Für durch Spannung $U$ beschleunigte Elektronen gilt der Energiesatz: $e \cdot U = \frac{1}{2} m_e v^2$
\end{tcolorbox}

\vspace{0.5em}

% ============================================================
% AUFGABE 2
% ============================================================

\textbf{Aufgabe 2: Beschleunigtes Elektron} \hfill \textbf{(6 BE)}

Ein Elektron wird aus der Ruhe durch die Spannung $U = \SI{150}{V}$ beschleunigt.

\begin{enumerate}[label=\alph*), leftmargin=*, nosep]
\item \textbf{Berechne} die Geschwindigkeit $v$ des Elektrons nach der Beschleunigung. \hfill (2 BE)

\vspace{0.3em}
$v = \sqrt{\dfrac{2 \cdot e \cdot U}{m_e}} = \sqrt{\dfrac{2 \cdot \lsg{$\SI{1,6e-19}{C}$} \cdot \lsg{$\SI{150}{V}$}}{\lsg{$\SI{9,1e-31}{kg}$}}}$

\vspace{0.5em}
$v = \sqrt{\lsg{$\SI{5,27e13}{m^2/s^2}$}} = $ \lsg{$\SI{7,3e6}{m/s}$}

\vspace{0.6em}

\item \textbf{Berechne} die De-Broglie-Wellenlänge $\lambda$ des Elektrons. \hfill (2 BE)

\vspace{0.3em}
$\lambda = \dfrac{h}{m_e \cdot v} = \dfrac{\lsg{$\SI{6,63e-34}{J.s}$}}{\lsg{$\SI{9,1e-31}{kg}$} \cdot \lsg{$\SI{7,3e6}{m/s}$}}$

\vspace{0.5em}
$\lambda = $ \lsg{$\SI{1,0e-10}{m}$} $=$ \lsg{100} pm

\vspace{0.6em}

\item \textbf{Vergleiche} $\lambda$ mit einem typischen Atomabstand in Kristallen ($d \approx \SI{200}{pm}$). Ist Elektronenbeugung an Kristallen zu erwarten? \textbf{Begründe.} \hfill (2 BE)

\vspace{0.3em}
\lsg{Die De-Broglie-Wellenlänge ($\lambda = \SI{100}{pm}$) ist von der gleichen Größenordnung wie Atomabstände in Kristallen ($\approx \SI{200}{pm}$). Daher ist Beugung an Kristallgittern möglich und beobachtbar (Davisson-Germer-Experiment 1927).}

\end{enumerate}

% ============================================================
% AUFGABE 3
% ============================================================

\textbf{Aufgabe 3: Materiewellen im Alltag?} \hfill \textbf{(7 BE)}

Ein Tennisball ($m = \SI{58}{g}$) wird mit $v = \SI{50}{m/s}$ geschlagen.

\begin{enumerate}[label=\alph*), leftmargin=*, nosep]
\item \textbf{Berechne} die De-Broglie-Wellenlänge des Tennisballs. \hfill (3 BE)

\vspace{0.3em}
$p = m \cdot v = $ \lsg{$\SI{0,058}{kg}$} $\cdot$ \lsg{$\SI{50}{m/s}$} $=$ \lsg{$\SI{2,9}{kg.m/s}$}

\vspace{0.3em}
$\lambda = \dfrac{h}{p} = \dfrac{\lsg{$\SI{6,63e-34}{J.s}$}}{\lsg{$\SI{2,9}{kg.m/s}$}} = $ \lsg{$\SI{2,3e-34}{m}$}

\vspace{0.6em}

\item Die kleinsten experimentell zugänglichen Strukturen (Atomgitter) haben Abstände von etwa $d \approx \SI{e-10}{m}$.

\textbf{Begründe}, warum wir bei Tennisbällen keine Beugungserscheinungen beobachten können. \hfill (4 BE)

\vspace{0.3em}
\lsg{Die De-Broglie-Wellenlänge des Tennisballs ($\lambda \approx \SI{2e-34}{m}$) ist um den Faktor $10^{24}$ kleiner als die kleinsten messbaren Strukturen ($\SI{e-10}{m}$). Um Beugung zu beobachten, müsste das Hindernis in der Größenordnung der Wellenlänge sein. Ein solches Hindernis existiert nicht -- es wäre etwa $10^{19}$-mal kleiner als ein Proton. Daher sind Welleneffekte bei makroskopischen Objekten prinzipiell nicht beobachtbar.}

\end{enumerate}

% ============================================================
% KASTEN 3: MERKSATZ
% ============================================================

\begin{tcolorbox}[merkbox={Kasten 3: Merksatz De-Broglie-Wellenlänge}]
Jedes bewegte Teilchen hat eine Wellenlänge \quad $\lambda = $ \lsg{$h$} $/$ \lsg{$p$}.

\vspace{0.3em}
Je größer die \lsg{Masse / der Impuls}, desto kleiner ist $\lambda$.

\vspace{0.3em}
Bei makroskopischen Objekten ist $\lambda$ so winzig (z.\,B.\ $10^{-34}$ m), dass

\vspace{0.3em}
\lsg{keine Beugungserscheinungen beobachtbar sind / Welleneffekte nicht messbar sind}.
\end{tcolorbox}

\vspace{0.3em}

\begin{center}
\rule{8cm}{0.5pt}

\textbf{Pflichtaufgaben:} \lsg{24} / 24 BE
\end{center}

\vspace{0.5em}

% ============================================================
% ZUSATZAUFGABEN
% ============================================================

{\large\textbf{Zusatzaufgaben}}

\vspace{0.3em}

\textbf{Z1:} \textbf{Berechne} die De-Broglie-Wellenlänge für ein Proton ($m_p = \SI{1,67e-27}{kg}$), das mit $v = \SI{5e6}{m/s}$ fliegt. \hfill (2 BE)

\vspace{0.3em}
$\lambda = $ \lsg{$\dfrac{h}{m_p \cdot v} = \dfrac{\SI{6,63e-34}{J.s}}{\SI{1,67e-27}{kg} \cdot \SI{5e6}{m/s}} = \SI{7,9e-14}{m} = \SI{0,079}{pm}$}

\vspace{0.5em}

\textbf{Z2:} Auf welche Spannung $U$ muss ein Elektron beschleunigt werden, damit es die Wellenlänge $\lambda = \SI{100}{pm}$ hat? \hfill (3 BE)

\vspace{0.3em}
\lsg{Aus $\lambda = \dfrac{h}{m_e \cdot v}$ folgt $v = \dfrac{h}{m_e \cdot \lambda} = \dfrac{\SI{6,63e-34}{J.s}}{\SI{9,1e-31}{kg} \cdot \SI{e-10}{m}} = \SI{7,3e6}{m/s}$}

\vspace{0.3em}
\lsg{Aus $v = \sqrt{\dfrac{2eU}{m_e}}$ folgt $U = \dfrac{m_e \cdot v^2}{2e} = \dfrac{\SI{9,1e-31}{kg} \cdot (\SI{7,3e6}{m/s})^2}{2 \cdot \SI{1,6e-19}{C}}$}

\vspace{0.3em}
$U = $ \lsg{$\SI{151}{V}$}

\vspace{0.5em}

\textbf{Z3:} Ein Elektron und ein Proton haben die gleiche kinetische Energie. Welches Teilchen hat die größere De-Broglie-Wellenlänge? \textbf{Begründe} mit einer Rechnung oder Argumentation. \hfill (3 BE)

\vspace{0.3em}
\lsg{Bei gleicher kinetischer Energie $E_{\text{kin}} = \frac{p^2}{2m}$ gilt $p = \sqrt{2m \cdot E_{\text{kin}}}$.}

\vspace{0.3em}
\lsg{Das schwerere Proton hat den größeren Impuls $p$, da $p \propto \sqrt{m}$.}

\vspace{0.3em}
\lsg{Da $\lambda = h/p$ gilt: größerer Impuls → kleinere Wellenlänge.}

\vspace{0.3em}
\lsg{\textbf{Das Elektron hat die größere De-Broglie-Wellenlänge} (um Faktor $\sqrt{m_p/m_e} \approx 43$).}

\vspace{0.5em}

\begin{center}
\rule{8cm}{0.5pt}

\textbf{Zusatzaufgaben:} \lsg{8} / 8 BE \hspace{2cm} \textbf{Gesamt:} \lsg{32} / 32 BE
\end{center}

\end{document}
